\chapter{Swaps de taux}\label{ch:m2:interest-rate-swaps}

Une entreprise a emprunté 200 millions de dollars pour dix ans au taux au jour le jour
plus 1.5~\%, et son conseil d'administration veut savoir combien les intérêts coûteront l'année prochaine
et chaque année suivante. Son trésorier appelle une banque, et celle-ci cote un
taux fixe de 4.05~\% sur dix ans~: l'entreprise paiera ce taux, la banque versera
à l'entreprise le taux au jour le jour, et l'intérêt variable du prêt est transformé
en un taux fixe de 5.55~\%. Le trader de la banque, qui vient de prendre l'autre côté
de dix ans de risque de taux, le couvre en moins d'une minute, sur le
marché interbancaire ou sur une plateforme électronique, au moyen d'un swap identique
conclu auprès d'un autre dealer ou de contrats à terme sur Treasury. Les dérivés de taux représentent
environ quatre cinquièmes du notionnel de l'ensemble des dérivés de gré à gré, et la
courbe des swaps est une référence par rapport à laquelle sont évalués et couverts les prêts bancaires, les obligations d'entreprises et
les produits structurés. Ce chapitre
évalue un swap, construit la courbe sur laquelle il est évalué, et en mesure le risque.

\section{Le swap et ses jambes}

\begin{definition}[Swap de taux, jambes fixe et variable]\label{def:m2:interest-rate-swaps:swap}
Un \emph{swap de taux}\index{swap de taux} (interest-rate swap) est un accord par lequel les parties s'engagent à
échanger, sur un montant notionnel qui n'est jamais échangé lui-même, des paiements
d'intérêts à taux fixe, la \emph{jambe fixe}\index{jambe fixe}, contre des paiements
à taux variable fixés par un taux de référence, la \emph{jambe variable}\index{jambe variable},
selon un calendrier de dates jusqu'à l'échéance.
\end{definition}

\begin{definition}[Swaps payeurs et receveurs]\label{def:m2:interest-rate-swaps:payer}
Dans un \emph{swap payeur}\index{swap payeur}, on paie le taux fixe et l'on reçoit
le taux variable~; dans un \emph{swap receveur}\index{swap receveur}, on reçoit le taux fixe et
paie le taux variable. Un payeur gagne lorsque les taux montent~: il est court en duration, comme
un emprunteur qui a fixé son taux, ou comme une position vendeuse sur une obligation.
\end{definition}

\begin{omfigure}
\begin{tikzpicture}[font=\small,
  box/.style={draw,thick,rounded corners=2pt,align=center,minimum height=1cm,minimum width=3cm,inner sep=4pt},
  arr/.style={-{Stealth[length=2.5mm]},thick}]
\node[box,draw=omDef,fill=omDef!8] (co) at (0,0) {entreprise\\(payeur)};
\node[box,draw=omThm,fill=omThm!8] (bk) at (6.4,0) {banque\\(receveur)};
\node[box,draw=gray,fill=gray!8] (ln) at (-5.2,0) {prêteurs};
\draw[arr] (co.20) -- node[above=5pt] {taux fixe de 4.05~\% par an} (bk.160);
\draw[arr] (bk.200) -- node[below=5pt] {SOFR composé} (co.340);
\draw[arr] (co) -- node[above=11pt,align=center] {SOFR + 1.50~\%} (ln);
\node[font=\footnotesize,text=omThm,align=center] at (3.2,-1.75) {notionnel de USD~200 millions~: jamais échangé,\\seulement utilisé pour calculer les paiements};
\end{tikzpicture}

\omcaption{La couverture de l'entreprise. Elle continue de verser à ses prêteurs le SOFR plus 1.50~\%~;
le swap lui verse le SOFR et lui facture 4.05~\%. Les paiements variables se compensent, et l'entreprise
paie un taux fixe de 5.55~\%. Chaque année, les paiements sont compensés et réglés deux
jours ouvrés après la fin de la période.}\label{fig:m2:interest-rate-swaps:flows}
\end{omfigure}

\begin{example}[Une année du swap de l'entreprise]\label{ex:m2:interest-rate-swaps:period}
Le swap prend effet le 29 septembre 2026. Le 29 septembre 2027, l'entreprise verse
$200\,000\,000 \times 4.05\% \times 365/360 = \text{USD}~8\,212\,500$ et reçoit le
taux au jour le jour composé sur ces 365 jours, quel qu'il soit~; sur
son prêt, l'entreprise verse le même taux composé, plus 1.5~\%. Les paiements variables
se compensent, et son coût est fixe.
\end{example}

\section{Swaps de taux au jour le jour et conventions}

\begin{definition}[Swap de taux au jour le jour]\label{def:m2:interest-rate-swaps:ois}
Un \emph{swap de taux au jour le jour}\index{swap de taux au jour le jour} (overnight index swap, OIS) est un swap
dont la \omterm{def:m2:interest-rate-swaps:swap}{jambe variable} verse un \omterm{def:m2:central-banks-and-the-short-rate:rfr}{taux de référence au jour le jour} \omterm{def:m2:central-banks-and-the-short-rate:arrears}{composé a posteriori} sur
chaque période (voir la \cref{def:m2:central-banks-and-the-short-rate:arrears}). Depuis la
fin du LIBOR, c'est le swap standard dans les devises dont le taux interbancaire
a cessé d'être publié~; en euros, il se négocie aux côtés des swaps sur Euribor, lequel
est toujours publié (voir le \cref{dat:m2:short-term-interest-rate-futures:contracts}).
\end{definition}

\begin{dated}{2026-09}{Conventions d'un swap SOFR en dollars}\label{dat:m2:interest-rate-swaps:conventions}
Les deux jambes paient annuellement, en actual/360~; la \omterm{def:m2:interest-rate-swaps:swap}{jambe variable} verse le SOFR \omterm{def:m2:central-banks-and-the-short-rate:arrears}{composé
a posteriori}, avec un paiement deux jours ouvrés du marché des titres d'État après la fin de
chaque période, de sorte que le dernier fixing soit connu~; le swap commence à courir deux
jours ouvrés après la date de transaction. Les conventions diffèrent selon la devise, et selon
l'échéance du swap (les swaps courts peuvent ne payer qu'une fois, à l'échéance).
\end{dated}

\section{Évaluation, taux au pair et risque}

\begin{definition}[Annuité et taux de swap au pair]\label{def:m2:interest-rate-swaps:par}
Pour les dates de paiement de la \omterm{def:m2:interest-rate-swaps:swap}{jambe fixe} $t_1, \dots, t_n$, avec les fractions de calcul $\delta_i$
et les facteurs d'actualisation $P(t_i)$, l'\emph{annuité}\index{annuité} est $A =
\sum_i \delta_i P(t_i)$, c'est-à-dire la valeur de 1 reçu chaque année selon le calendrier
de la \omterm{def:m2:interest-rate-swaps:swap}{jambe fixe}. Le \emph{taux de swap au pair}\index{taux de swap au pair} est le taux fixe auquel le swap
vaut zéro à sa date de départ.
\end{definition}

\begin{proposition}[Valeur d'un OIS sur une courbe unique]\label{prop:m2:interest-rate-swaps:value}
Si les taux au jour le jour sont actualisés sur la même courbe que celle qui les projette, la
\omterm{def:m2:interest-rate-swaps:swap}{jambe variable} d'un OIS entre $t_0$ et $t_n$ vaut $P(t_0) - P(t_n)$ par unité
de notionnel, quel que soit son calendrier. Ainsi, la valeur du payeur est
\[
V \;=\; N\bigl[P(t_0) - P(t_n) - K\,A\bigr],
\qquad
K_{\text{pair}} \;=\; \frac{P(t_0) - P(t_n)}{A},
\]
et $\partial V/\partial K = -N A$~: la sensibilité du swap à son propre taux est
son \omterm{def:m2:interest-rate-swaps:par}{annuité}.
\end{proposition}
\begin{proof}
Le paiement composé au jour le jour d'une période en $t_i$ se réplique en investissant
$P(t_{i-1})$ en $t_0$ dans un dépôt jusqu'à $t_{i-1}$, puis en reconduisant 1 au jour le jour
à partir de là, ce qui donne $1/P(t_{i-1}, t_i)$ en $t_i$, dont 1 est restitué et le
reste est le paiement~; sa valeur vaut $P(t_{i-1}) - P(t_i)$, et la somme se télescope.
\end{proof}

\begin{method}[Bootstrap d'une courbe OIS]\label{met:m2:interest-rate-swaps:bootstrap}
Prendre les \omterm{def:m2:interest-rate-swaps:par}{taux de swap au pair} pour des échéances croissantes. Résoudre pour le facteur d'actualisation
à la première échéance pour que son swap vaille zéro~; puis, en le conservant, pour la
deuxième, en interpolant les facteurs d'actualisation de toutes les dates de paiement intermédiaires
entre les pivots (ici log-linéairement en temps)~; et ainsi de suite. Vérifier que chaque
swap d'entrée se réévalue bien au pair.
\end{method}

\begin{example}[Un swap à dix ans]\label{ex:m2:interest-rate-swaps:ten}
Avec des \omterm{def:m2:government-bonds:zero}{taux au pair} illustratifs de 3.90, 3.85, 3.83, 3.85, 3.92 et 4.05~\% à un,
deux, trois, cinq, sept et dix ans, le \omterm{def:m2:government-bonds:zero}{taux au pair} à dix ans vaut 4.05~\% par
construction, l'\omterm{def:m2:interest-rate-swaps:par}{annuité} vaut 8.228 et le DV01 du swap sur 100~millions d'USD
vaut 82\,278~USD, soit $N A \times 10^{-4}$, et 82\,262~USD lorsque la courbe est reconstruite après
avoir décalé chaque donnée d'entrée d'un point de base. Tout cela se concentre sur le pivot à dix ans~:
un swap au pair est couvert exactement par le swap qui a construit son pivot.
\end{example}

\begin{omfigure}
\begin{tikzpicture}
\begin{axis}[width=11cm,height=5cm,
  xlabel={échéance (années)},ylabel={taux (\%)},
  tick label style={font=\small},label style={font=\small},
  xmin=0,xmax=10.5,ymin=3.7,ymax=4.5,grid=major,grid style={gray!25},
  legend columns=3,legend style={font=\footnotesize,at={(0.5,-0.3)},anchor=north,draw=none,fill=none,/tikz/every even column/.append style={column sep=8pt}},
  table/col sep=comma]
\addplot[only marks,mark=*,mark size=2pt,omDef] table[x=years,y=par]{figdata/markets-2/09-interest-rate-swaps/par.csv};
\addplot[omThm,very thick] table[x=years,y=zero]{figdata/markets-2/09-interest-rate-swaps/curve.csv};
\addplot[omProp,thick,const plot] table[x=years,y=fwd]{figdata/markets-2/09-interest-rate-swaps/curve.csv};
\legend{taux de swap au pair (données d'entrée),taux zéro,taux forward à un an}
\end{axis}
\end{tikzpicture}

\omcaption{La courbe obtenue par bootstrap à partir de six \omterm{def:m2:government-bonds:zero}{taux au pair} illustratifs. Les taux zéro
sont réguliers~; les taux forward sont plats entre les pivots et sautent à chacun d'eux, un
artefact de l'interpolation linéaire du logarithme des facteurs d'actualisation~: tout
choix d'interpolation est un choix de courbe forward, et le livre~6
traite de meilleurs choix. Données~: le tutoriel du chapitre.}\label{fig:m2:interest-rate-swaps:curve}
\end{omfigure}

\begin{omfigure}
\begin{tikzpicture}
\begin{axis}[width=11cm,height=4.8cm,ybar,bar width=12pt,
  ylabel={DV01 (milliers d'USD par pb)},
  tick label style={font=\small},label style={font=\small},
  xmin=-0.6,xmax=5.6,ymin=-55,ymax=95,grid=major,grid style={gray!25},
  xtick=data,xticklabels={1 an,2 ans,3 ans,5 ans,7 ans,10 ans},
  legend columns=2,legend style={font=\footnotesize,at={(0.5,-0.22)},anchor=north,draw=none,fill=none,/tikz/every even column/.append style={column sep=8pt}},
  table/col sep=comma]
\addplot[fill=omDef!55,draw=omDef] table[x=i,y=par10]{figdata/markets-2/09-interest-rate-swaps/buckets.csv};
\addplot[fill=omThm!55,draw=omThm] table[x=i,y=fwd5y5y]{figdata/markets-2/09-interest-rate-swaps/buckets.csv};
\legend{swap payeur au pair à 10 ans,swap payeur forward 5y5y}
\end{axis}
\end{tikzpicture}

\omcaption{DV01 par compartiment de deux \omterm{def:m2:interest-rate-swaps:payer}{swaps payeurs} de 100~millions d'USD~: la variation
de valeur lorsque chaque \omterm{def:m2:government-bonds:zero}{taux au pair} d'entrée augmente d'un point de base et que la courbe est reconstruite.
Le swap au pair à dix ans ne charge que son propre pivot. Le swap de cinq ans débutant dans cinq
ans est acheteur des taux à dix ans et vendeur des taux à cinq ans~: c'est un pari sur le
forward, et sa couverture est une paire de swaps au pair. Données~: le
tutoriel du chapitre.}\label{fig:m2:interest-rate-swaps:buckets}
\end{omfigure}

\section{Spreads de swap}

\begin{definition}[Spread de swap]\label{def:m2:interest-rate-swaps:spread}
Le \emph{spread de swap}\index{spread de swap} à une échéance donnée est le \omterm{def:m2:interest-rate-swaps:par}{taux de swap au pair}
moins le rendement de l'obligation d'État de même échéance.
\end{definition}

Pendant des décennies, les taux de swap étaient supérieurs aux rendements des obligations d'État, car la \omterm{def:m2:interest-rate-swaps:swap}{jambe variable}
était un taux interbancaire portant un risque de crédit bancaire. En septembre 2008, peu
après la faillite de Lehman Brothers, le \omterm{def:m2:interest-rate-swaps:spread}{spread de swap} en dollars à trente ans passa
en territoire négatif, et il y est resté depuis. Un \omterm{def:m2:interest-rate-swaps:ois}{swap de taux au jour le jour} porte
presque aucun risque de crédit dans sa \omterm{def:m2:interest-rate-swaps:swap}{jambe variable}, donc un spread négatif dit
autre chose~: détenir un \omterm{def:m2:money-markets:bill}{bon du Trésor} américain mobilise un bilan à financer au repo,
tandis que recevoir le fixe sur un swap ne demande que de la marge, et la demande de long terme pour
la duration, émanant des fonds de pension qui couvrent leurs passifs, rencontre des dealers dont
les bilans sont coûteux. Le \omterm{def:m2:interest-rate-swaps:spread}{spread de swap} est le prix de cet écart,
et le négocier, en étant acheteur de \omterm{def:m2:money-markets:bill}{bons du Trésor} contre paiement du fixe ou l'inverse, est
l'un des principaux arbitrages de valeur relative du marché des taux (livre~9).

\section{Plateformes d'exécution et compression}

\begin{definition}[Swap execution facility et compression de portefeuille]\label{def:m2:interest-rate-swaps:sef}
Une \emph{swap execution facility}\index{swap execution facility} (SEF) est une
plateforme de négociation réglementée pour les swaps aux États-Unis~; les swaps ayant
été déclarés négociables doivent être exécutés sur l'une d'elles. \emph{Compression
de portefeuille}\index{compression de portefeuille} consiste à résilier des swaps compensatoires
entre plusieurs contreparties, en les remplaçant par moins de swaps présentant le même
risque net, afin de réduire le notionnel brut, le nombre de transactions et le capital
et les marges qu'ils mobilisent.
\end{definition}

Le marché des swaps était, jusqu'aux règles adoptées après la crise de 2008, largement un
marché bilatéral entre dealers et leurs clients, négocié par téléphone. Il repose désormais sur trois niveaux~:
compensation centrale pour la plupart des swaps standard (voir le \cref{ch:m2:swap-clearing-and-the-ccp-basis}),
exécution obligatoire des plus standardisés sur des plateformes électroniques, et
déclaration de chaque transaction. La compression vient en complément~: comme les swaps ne sont jamais
dénoués par revente, mais seulement compensés par de nouveaux swaps, le notionnel brut croît
à chaque couverture, sauf s'il est périodiquement résilié.

\begin{omfigure}
\begin{tikzpicture}[font=\small,
  node/.style={draw,thick,circle,minimum size=1.1cm,align=center},
  arr/.style={-{Stealth[length=2.5mm]},thick}]
\node[node,draw=omDef,fill=omDef!8] (a) at (0,1.6) {A};
\node[node,draw=omDef,fill=omDef!8] (b) at (2.2,-0.9) {B};
\node[node,draw=omDef,fill=omDef!8] (c) at (-2.2,-0.9) {C};
\draw[arr,omThm] (a) -- node[right=6pt,align=left] {A paie le fixe\\à B} (b);
\draw[arr,omThm] (b) -- node[below=4pt] {B paie le fixe à C} (c);
\draw[arr,omThm] (c) -- node[left=6pt,align=right] {C paie le fixe\\à A} (a);
\node[align=center,anchor=west] at (4.8,0.4) {notionnel brut avant~: 300\\risque net de chaque dealer~: 0\\après compression~: rien};
\end{tikzpicture}

\omcaption{Un cycle de compression. Trois dealers détiennent trois swaps identiques en
cercle~: chacun paie le fixe sur l'un et reçoit sur un autre, si bien qu'aucun n'a de risque, pourtant
tous trois immobilisent un notionnel brut, des marges et du capital pour les couvrir. La compression
résilie le cycle. Les cycles réels impliquent de nombreux participants et des swaps qui
ne se compensent qu'approximativement~; le service identifie le plus grand ensemble de
résiliations laissant le risque de chaque participant dans ses
seuils de tolérance.}\label{fig:m2:interest-rate-swaps:compression}
\end{omfigure}

\begin{dated}{2026-09}{La taille du marché}\label{dat:m2:interest-rate-swaps:market}
Notionnel des produits dérivés de gré à gré en cours~: 846~billions de dollars à fin juin
2025 (la dernière enquête semestrielle de la BRI), soit 16~\% de plus qu'un an plus tôt~;
les produits dérivés de taux représentaient 79~\% de ce total, soit environ 668~billions de dollars,
les contrats en euros dépassant les contrats en dollars depuis 2022. Aux États-Unis, l'inscription
auprès d'un \omterm{def:m2:interest-rate-swaps:sef}{swap execution facility} (SEF) était obligatoire depuis octobre 2013, et
l'exécution sur un SEF des swaps déclarés disponibles à la négociation à compter du 15 février 2014.
\end{dated}

\section{Tutoriel~: une courbe et un swap}\label{tut:m2:interest-rate-swaps:curve}

\begin{tutorial}
\textbf{Objectif.} Construire par bootstrap une courbe OIS à partir de six \omterm{def:m2:government-bonds:zero}{taux au pair},
recalculer le prix des instruments d'entrée, et calculer le DV01 segmenté d'un swap au pair et d'un swap à départ différé.
\textbf{État final~:}
Les \cref{fig:m2:interest-rate-swaps:curve,fig:m2:interest-rate-swaps:buckets} et les
résultats de l'\cref{ex:m2:interest-rate-swaps:ten}.
\begin{tutsteps}
\item \textbf{\omterm{def:m2:interest-rate-swaps:par}{Annuité}, \omterm{def:m2:government-bonds:zero}{taux au pair}, valeur} (la \cref{prop:m2:interest-rate-swaps:value}).
\omcode{code/firm/curve/firm_curve.py}{50}{63}{Annuité, taux au pair et valeur d'un OIS dans une courbe unique.}\label{lst:m2:interest-rate-swaps:value}
\item \textbf{Bootstrap}~: un pilier à la fois, par dichotomie sur le logarithme
du facteur d'actualisation.
\omcode{code/firm/curve/firm_curve.py}{66}{84}{Bootstrap des facteurs d'actualisation de sorte que chaque swap d'entrée soit au pair.}\label{lst:m2:interest-rate-swaps:bootstrap}
Chaque swap d'entrée doit se réévaluer à zéro, à un millième de dollar près.
\item \textbf{DV01 segmenté}~: appliquer un choc à chaque entrée, reconstruire la courbe, puis réévaluer le prix.
\omcode{code/firm/curve/firm_curve.py}{87}{95}{DV01 segmenté par choc et reconstruction.}\label{lst:m2:interest-rate-swaps:buckets}
\item \textbf{Exécuter} \texttt{swaps\_demo.ten\_year()} et
\texttt{swaps\_demo.forward\_5y5y()}~; \texttt{fig\_swaps.py} produit les graphiques.
\end{tutsteps}
\textbf{Pour aller plus loin.} Remplacer l'interpolation log-linéaire des facteurs d'actualisation
par une interpolation linéaire des taux zéro, et observer la courbe forward se modifier
alors que chaque entrée se réévalue toujours au pair~; puis évaluer un swap dont le taux fixe
dépasse de 1~\% le \omterm{def:m2:government-bonds:zero}{taux au pair}, et lire ses segments.
\end{tutorial}

\section{Construction~: la courbe OIS et le valorisateur de swaps}\label{bld:m2:interest-rate-swaps:curve}

\begin{build}
\bfield{Rôle} Chaque position en taux de la mini-entreprise, et chaque
actualisation d'un flux futur lit cette courbe. C'est la première courbe de la
bibliothèque de valorisation (les livres~5 et~6 l'étendent à plusieurs courbes, aux
\omterm{def:m2:money-markets:turn}{passages de fin de période} et aux dates de réunion, voir la \cref{bld:m2:short-term-interest-rate-futures:meetings}).
\bfield{Interface} \texttt{Curve(spot, times, dfs)} avec \texttt{df(date)},
\texttt{zero(date)}~; \texttt{schedule(start, years)}~; \texttt{annuity}~;
\texttt{par\_rate}~; \texttt{swap\_pv(curve, dates, fixed, notional, payer)}~;
\texttt{bootstrap(spot, tenors, rates)}~; \texttt{bucket\_dv01(\ldots)}.
\bfield{Règles} Périodes annuelles, calcul des intérêts en actual/360, sans ajustement
des jours ouvrés dans cette version~; interpolation log-linéaire des facteurs d'actualisation,
taux forward plat au-delà du dernier pilier~; une seule courbe pour la projection et l'actualisation.
\bfield{Tests d'acceptation} \texttt{code/firm/curve/tests/}~: chaque entrée se réévalue au pair~;
une courbe d'entrée plate donne des \omterm{def:m2:government-bonds:zero}{taux au pair} constants aux échéances intermédiaires~;
un payeur gagne lorsque les taux montent~; un swap au pair ne charge que son propre pilier~;
un forward 5y5y est vendeur du pilier à cinq ans et acheteur du pilier à dix ans~; jours bissextiles.
\bfield{Pour aller plus loin} Calendriers de jours ouvrés et délais de paiement (les conventions de
la \cref{dat:m2:interest-rate-swaps:conventions})~; contrats à terme et dates de réunion à court
terme~; une courbe d'actualisation distincte, traitée au \cref{ch:m2:swap-clearing-and-the-ccp-basis}.
\end{build}

\begin{omsources}
\item Bank for International Settlements, \emph{OTC derivatives statistics at
end-June 2025}, décembre 2025.
\item S.\ Klingler and S.\ Sundaresan, «~An explanation of negative swap spreads:
demand for duration from underfunded pension plans~», \emph{Journal of
Finance} 74, 2019.
\item Clarus Financial Technology, «~SOFR swap nuances~»~; ARRC, \emph{An
Updated User's Guide to SOFR}, 2021.
\item CFTC, règle finale sur les principes fondamentaux et les autres exigences applicables aux swap
execution facilities, 2013.
\end{omsources}

\section{Exercices}

\begin{exercise}[$\star$]\label{exo:m2:interest-rate-swaps:1}
Pour un \omterm{def:m2:interest-rate-swaps:payer}{swap payeur} de 100~millions de dollars à 4.05~\%, dont la période annuelle compte 365 jours,
calculer le paiement fixe dû à l'échéance et dire ce que reçoit le payeur.
\end{exercise}

\begin{exercise}[$\star$]\label{exo:m2:interest-rate-swaps:2}
Un fonds de pension reçoit le fixe sur un swap trentenaire. Déterminer si sa duration est
longue ou courte, et dire s'il gagne ou perd si les taux montent.
\end{exercise}

\begin{exercise}[$\star$]\label{exo:m2:interest-rate-swaps:3}
Le taux de swap à dix ans est de 4.05~\% et le Treasury à dix ans rapporte 4.20~\%. Calculer le \omterm{def:m2:interest-rate-swaps:spread}{spread de swap} et
dire ce qu'il aurait signifié avant 2008.
\end{exercise}

\begin{exercise}[$\star\star$]\label{exo:m2:interest-rate-swaps:4}
Le DV01 du swap à dix ans est de 82\,278~\$ à partir de l'\omterm{def:m2:interest-rate-swaps:par}{annuité} et de 82\,262~\$ en appliquant un choc
puis en reconstruisant la courbe. Expliquer l'écart, et dire lequel retenir pour la couverture.
\end{exercise}

\begin{exercise}[$\star\star$]\label{exo:m2:interest-rate-swaps:5}
Le swap forward 5y5y de la \cref{fig:m2:interest-rate-swaps:buckets}, dont le \omterm{def:m2:government-bonds:zero}{taux au pair} est de 4.295~\%,
expliquer son DV01 segmenté et donner les swaps au pair qui le couvrent.
\end{exercise}

\begin{exercise}[$\star\star$]\label{exo:m2:interest-rate-swaps:6}
Le dealer A paie le fixe à B, B paie le fixe à C et C paie le fixe à A, chacun pour
100~millions de dollars d'un même swap à dix ans. Calculer le notionnel brut et le risque net de chaque dealer,
puis le résultat de la \omterm{def:m2:interest-rate-swaps:sef}{compression de portefeuille}.
\end{exercise}

\begin{exercise}[$\star\star\star$]\label{exo:m2:interest-rate-swaps:7}
\emph{Codage.} Avec \texttt{firm\_curve}, construire par bootstrap la courbe de
l'\cref{ex:m2:interest-rate-swaps:ten} et donner les \omterm{def:m2:government-bonds:zero}{taux au pair} à quatre ans et à huit ans
ainsi que le facteur d'actualisation à dix ans.
\end{exercise}

\begin{exercise}[$\star\star\star$]\label{exo:m2:interest-rate-swaps:8}
\emph{Trouver l'erreur.} «~Le \omterm{def:m2:interest-rate-swaps:spread}{spread de swap} trentenaire est négatif, donc le marché juge que les
banques sont plus sûres que le gouvernement américain.~» Corriger cette affirmation.
\end{exercise}

\section{Problème~: La couverture d'entreprise}

\begin{problem}[{Problème du week-end --- fixation d'un emprunt à dix ans}]\label{pb:m2:interest-rate-swaps:1}
Une entreprise a emprunté 200~millions de dollars sur dix ans, au taux SOFR \omterm{def:m2:central-banks-and-the-short-rate:arrears}{composé a posteriori}
majoré de 1.50~\%. Le 25 septembre 2026, elle entre dans un \omterm{def:m2:interest-rate-swaps:payer}{swap payeur} de dix ans débutant le 29
septembre au \omterm{def:m2:interest-rate-swaps:par}{taux de swap au pair} de la courbe de l'\cref{ex:m2:interest-rate-swaps:ten}.

\medskip\noindent\textbf{Partie~I --- Le swap.}
\begin{enumerate}
\item Donner le taux fixe et le coût fixe tout compris de l'entreprise.
\item Donner le DV01 du swap pour l'entreprise, avec son signe.
\item Quelle est la valeur du swap pour l'entreprise à la date de transaction~?
\item Quel pilier porte son risque~?
\item Pourquoi le swap débute-t-il le 29 septembre et non à la date de transaction~?
\end{enumerate}

\medskip\noindent\textbf{Partie~II --- Scénarios.}
\begin{enumerate}[resume]
\item Tous les \omterm{def:m2:government-bonds:zero}{taux au pair} montent de 50 points de base. Quelle est la valeur du swap~?
\item Tous baissent de 50 points de base.
\item Pourquoi les deux valeurs ne sont-elles pas égales et opposées~?
\item La courbe se pentifie~: $-20$ points de base à un an, croissant linéairement jusqu'à $+15$
à dix ans. Donner les chocs à chaque pilier et la valeur du swap.
\item Pourquoi le steepener vaut-il presque exactement 15 fois le DV01~?
\end{enumerate}

\medskip\noindent\textbf{Partie~III --- Les comptes.}
\begin{enumerate}[resume]
\item Après la baisse de 50 points de base, le swap affiche une perte. L'entreprise a-t-elle
perdu de l'argent~?
\item Que cherche à obtenir la comptabilité de couverture, et pourquoi les trésoriers y tiennent-ils~?
\item La banque a coté 4.07~\% au lieu du milieu à 4.05~\%. Combien l'entreprise a-t-elle payé
pour la couverture, en dollars d'aujourd'hui~?
\item Quel risque de crédit chaque partie supporte-t-elle, et comment est-il géré~?
\item Que laisserait non couvert un swap à cinq ans~?
\end{enumerate}

\medskip\noindent\textbf{Partie~IV --- Jugement.}
\begin{enumerate}[resume]
\item Pourquoi une entreprise pourrait-elle préférer un swap à l'émission d'une obligation à taux fixe~?
\item Comment la banque couvre-t-elle le swap qu'elle vient de conclure~?
\item Pourquoi la couverture de l'entreprise est un \omterm{def:m2:interest-rate-swaps:payer}{swap payeur}, et celle d'un fonds de pension
un \omterm{def:m2:interest-rate-swaps:payer}{swap receveur}~?
\item Énoncer le \emph{résultat nommé}~: le taux fixe, le DV01 de la couverture et sa
valeur après la pentification.
\item En une phrase~: qu'échange un swap~?
\end{enumerate}
\end{problem}

\section{Questions d'entretien}

\begin{interviewq}[$\star$ \iqroles{trader, banque}]\label{iq:m2:interest-rate-swaps:1}
Qu'est-ce qu'un \omterm{def:m2:interest-rate-swaps:swap}{swap de taux}, et pourquoi une entreprise en conclurait-elle un~?
\end{interviewq}

\begin{interviewq}[$\star$ \iqroles{chercheur, banque}]\label{iq:m2:interest-rate-swaps:2}
Comment valorisez-vous la \omterm{def:m2:interest-rate-swaps:swap}{jambe variable} d'un \omterm{def:m2:interest-rate-swaps:ois}{swap de taux au jour le jour} (OIS)~?
\end{interviewq}

\begin{interviewq}[$\star\star$ \iqroles{chercheur, développeur}]\label{iq:m2:interest-rate-swaps:3}
Décrivez-moi la construction par bootstrap d'une courbe d'actualisation à partir des taux de swap. Quels
choix faites-vous, et comment testez-vous le résultat~?
\end{interviewq}

\begin{interviewq}[$\star\star$ \iqroles{trader, chercheur}]\label{iq:m2:interest-rate-swaps:4}
Pourquoi les \omterm{def:m2:interest-rate-swaps:spread}{spreads de swap} trentenaires américains sont-ils négatifs depuis 2008~?
\end{interviewq}

\begin{interviewq}[$\star\star$ \iqroles{trader}]\label{iq:m2:interest-rate-swaps:5}
Vous recevez le taux fixe d'un swap à terme 5y5y. Quel est votre risque, et comment
le couvrez-vous avec des swaps au pair~?
\end{interviewq}

\begin{interviewq}[$\star\star\star$ \iqroles{développeur, chercheur}]\label{iq:m2:interest-rate-swaps:6}
Vos DV01 par compartiment ne redonnent pas le DV01 parallèle. Énumérez les causes possibles.
\end{interviewq}
